\documentclass[12pt,a4paper]{article}
\usepackage[top=22mm,bottom=24mm,left=27mm,right=27mm]{geometry}
\usepackage{amsmath,amssymb,xcolor}
\definecolor{link}{HTML}{205973}
\usepackage[hypertexnames=false,colorlinks=true,linkcolor=link,urlcolor=link]{hyperref}
\hypersetup{pdftitle={Theorem 3. Uniform concentration for a fixed policy}}
\setlength{\parindent}{0pt}
\setlength{\parskip}{5pt}
\raggedbottom
\newcommand{\R}{\mathbb R}
\newcommand{\E}{\mathbb E}
\newcommand{\Prob}{\mathbb P}
\newcommand{\F}{\mathcal F}
\newcommand{\Cbar}{\overline C}
\newcommand{\resclass}{\overline{\mathcal F}_{\!\mathcal S,\mu}}
\newcommand{\lead}[1]{\par\smallskip\noindent\textbf{#1}\enspace}
\begin{document}
{\LARGE\bfseries Theorem 3. Uniform concentration\\[2pt] for a fixed policy\par}
\medskip
How much can a sample underestimate the expected positive Bellman residual?
Theorem~3 controls this error simultaneously for all admissible value
functions. A value function may therefore be chosen after observing the
sample. The policy and the distribution of the sampled states are fixed
before the draw.

\section*{The residual class and the guarantee}
Let \(\mathcal S\) and \(\mathcal A\) be finite nonempty state and action
sets, with cost \(g_{x,a}\), transition probabilities \(p(x,y,a)\), and
discount \(0<\alpha<1\). For any policy \(\mu:\mathcal S\to\mathcal A\), write
\[
 P_\mu(x,y)=p(x,y,\mu(x)),\quad g_\mu(x)=g_{x,\mu(x)},\quad
 T_\mu J=g_\mu+\alpha P_\mu J .
\]
For a feature map \(\Phi:\mathcal S\to\mathcal H\) into a real Hilbert space,
write \(\widetilde J_{\mathbf z,b}(x)=\langle\Phi(x),\mathbf z\rangle+b\).
For radii \(C,B\ge0\), define
\[
 \resclass=\bigl\{(\widetilde J_{\mathbf z,b}-T_\mu\widetilde J_{\mathbf z,b})_+
                 :\ \|\mathbf z\|\le C,\ |b|\le B\bigr\},\qquad u_+=\max(u,0).
\]
Each member measures how far the fit exceeds the immediate cost plus
the discounted expected continuation value under \(\mu\). Let \(X_1,\ldots,X_N\) be independent draws from any
probability distribution \(\pi\) on \(\mathcal S\). Write
\(P f=\sum_x\pi_x f(x)\), \(P_N f=N^{-1}\sum_{i=1}^N f(X_i)\), and set
\[
 \begin{aligned}
 K&=\max_x\|\Phi(x)\|,\qquad \|g\|_\infty=\max_{x,a}|g_{x,a}|,\\
 \Cbar&=4CK(1+\alpha)+4B(1-\alpha)+2\|g\|_\infty,\\
 \Xi(C,B,K,\delta)&=\Cbar\bigl(1+\sqrt{2\ln(1/\delta)}\bigr).
 \end{aligned}
\]

\textbf{Theorem 3.} Fix \(\epsilon>0\) and \(0<\delta<1\). For every integer
\(N\ge1\) satisfying \(N\ge \Xi(C,B,K,\delta)^2/\epsilon^2\),
\begin{equation}
 \Prob\!\left\{\sup_{f\in\resclass}(P f-P_N f)\ge\epsilon\right\}\le\delta .
 \tag{T3}\label{eq:theorem}
\end{equation}
In particular, equality at \(\epsilon\) is included in the failure event.
Taking \(\mu\) optimal and \(\pi\) its discounted occupancy gives the
uniform sampling estimate used in Theorems~1 and~2.

\lead{A bound on one residual.}
Define \(d_\mu(x)=\Phi(x)-\alpha\sum_yP_\mu(x,y)\Phi(y)\). The transition
weights sum to one, so \(\|d_\mu(x)\|\le K(1+\alpha)\). For the residual
\(f\) of \(\widetilde J_{\mathbf z,b}\) this gives
\begin{equation}
 \begin{aligned}
 f(x)&=\bigl(\langle d_\mu(x),\mathbf z\rangle+(1-\alpha)b-g_\mu(x)\bigr)_+,\\
 0\le f(x)&\le CK(1+\alpha)+(1-\alpha)B+\|g\|_\infty\le \Cbar/2 .
 \end{aligned}
 \tag{1}\label{eq:residual}
\end{equation}

\newpage
\section*{A one-sided bound for any bounded class}
For a nonempty class \(\F\) of real functions on \(\mathcal S\), let
\(\sigma_1,\ldots,\sigma_N\) be independent uniform signs in \(\{-1,1\}\),
independent of the states. The paper's Rademacher complexity is
\[
 \widehat R_N(\F)=\E_\sigma\!\left[\sup_{f\in\F}\frac{2}{N}
          \sum_i\sigma_i f(X_i)\right],\qquad
 R_N(\F)=\E_X\widehat R_N(\F).
\]
The signs are averaged with the sample held fixed in \(\widehat R_N\).

\textbf{Lemma 1 (positive envelope).} Let \(N\ge1\), \(0<\delta<1\), and
suppose \(|f(x)|\le B\) for every \(f,x\). Here \(B>0\) is any uniform bound on the class. Then
\begin{equation}
 Z=\sup_{f\in\F}(P f-P_N f),\qquad
 \Prob\!\left\{Z\ge R_N(\F)+B\sqrt{\frac{2\ln(1/\delta)}{N}}\right\}\le\delta .
 \tag{2}\label{eq:lemma}
\end{equation}
The positive envelope matters for this closed event. With \(B=0\) and
\(\F=\{0\}\), \(Z=R_N(\F)=0\) on every sample. The probability on the
left is then exactly one for every \(0<\delta<1\).

\lead{Arbitrary classes.}
Regard \(\F\) as a subset of \(\R^{\mathcal S}\) with its pointwise topology.
A countable base of rational rectangles gives a countable dense subclass
by choosing one member of \(\F\) from each basic open set it meets.
Evaluation and finite sums are continuous, so this subclass gives the
same suprema in \(Z\) and \(\widehat R_N\), including when no maximum exists.
On countably infinite discrete state spaces the same argument uses basic
sets restricting finitely many coordinates. There \(P f\) is continuous
on the bounded class by dominated convergence. The suprema are therefore
measurable for both finite and countable discrete state spaces.

\lead{Symmetrization.}
Draw an independent copy \(X'_1,\ldots,X'_N\) and write \(P'_N\) for
its empirical mean. Since \(\E_{X'}P'_N f=P f\), moving the supremum
inside this expectation gives \(\E Z\le\E\sup_f(P'_N f-P_N f)\).
Swapping the two draws of each pair independently preserves their joint law.
Define \(A(f,X,\sigma)=N^{-1}\sum_i\sigma_i f(X_i)\). The swap gives
\[
 \E Z\le\E\sup_f\bigl(A(f,X',\sigma)+ A(f,X,-\sigma)\bigr)
 \le \E\sup_f A(f,X',\sigma)+\E\sup_f A(f,X,-\sigma).
\]
Each term equals \(R_N(\F)/2\) because \(-\sigma\) has the same law as
\(\sigma\). This accounts for the factor \(2/N\).

\lead{Bounded differences.}
Replacing one draw changes every \(P_N f\), hence \(Z\), by at most
\(c=2B/N\). Reveal the draws successively. The conditional-mean increments
of \(Z\) are centered. Each has conditional range of width at most \(c\)
because averaging over future draws preserves the coordinate bound.
The second derivative of its log moment-generating function is a tilted
variance, bounded by the squared half-width \(c^2/4\). Integrating twice
and iterating conditional expectations yields
\[
 \E\exp(\lambda(Z-\E Z))\le\exp(N\lambda^2c^2/8).
\]
Markov's inequality at \(\lambda=4t/(Nc^2)\) gives
\(\Prob\{Z-\E Z\ge t\}\le\exp(-2t^2/(Nc^2))\) for \(t>0\).
Taking \(t=B\sqrt{2\ln(1/\delta)/N}\) and using \(\E Z\le R_N(\F)\)
proves \eqref{eq:lemma}.

\newpage
\section*{The complexity of a fixed-policy residual}
Fix the sample and abbreviate \(d_i=d_\mu(X_i)\). By \eqref{eq:residual},
each residual applies the 1-Lipschitz map
\(\varphi_i(u)=(u-g_\mu(X_i))_+\) to the affine score
\(u_i=\langle d_i,\mathbf z\rangle+(1-\alpha)b\).
For this one-sided complexity, such coordinatewise transformations can
only decrease the sign average.

To see this, condition on all signs except \(\sigma_i\). Index the candidates
by \(j\), write \(u_j\) for their scores at coordinate \(i\), and \(a_j\)
for their signed sums over the other coordinates. Twice the conditional
expectation of the unnormalized supremum is
\[
 \sup_{j,k}\{a_j+a_k+\varphi_i(u_j)-\varphi_i(u_k)\}
 \le \sup_{j,k}\{a_j+a_k+|u_j-u_k|\}.
\]
Exchanging \(j,k\) leaves \(a_j+a_k\) unchanged, so the last supremum equals
\(\sup_{j,k}(a_j+a_k+u_j-u_k)\). This is twice the conditional expectation
for the untransformed coordinate. Repeating for each \(i\) removes all
\(\varphi_i\) while preserving the inequality.

\lead{Feature weights and intercept.}
For fixed signs, Cauchy--Schwarz bounds the linear supremum by
\(C\|\sum_i\sigma_i d_i\|\). The corresponding intercept bound is
\((1-\alpha)B|\sum_i\sigma_i|\). Independence and zero means cancel the
cross terms in both second moments, giving
\[
 \E_\sigma\Big\|\sum_i\sigma_i d_i\Big\|^2=\sum_i\|d_i\|^2,
 \qquad \E_\sigma\Big|\sum_i\sigma_i\Big|^2=N .
\]
Cauchy--Schwarz and \(\|d_i\|\le K(1+\alpha)\) thus give, on every sample,
\begin{equation}
 \widehat R_N(\resclass)
 \le \frac{2}{\sqrt N}\bigl(CK(1+\alpha)+(1-\alpha)B\bigr)
 \le \frac{\Cbar}{\sqrt N}.
 \tag{3}\label{eq:complexity}
\end{equation}
Averaging proves the same bound for \(R_N\). The cost enters the pointwise
bound \eqref{eq:residual} through \(\|g\|_\infty\). The translation by \(-g_\mu(X_i)\)
preserves the Lipschitz constant.

\lead{The stated sample size.}
If \(\Cbar>0\), use \eqref{eq:residual} to set the generic bound of Lemma~1 to
\(\Cbar/2\). The intercept radius \(B\) still enters through \(\Cbar\).
The final display on page~30 of the manuscript uses the unbarred linear
class \(\mathcal F_{\mathcal S,\mu}\) and the symbol \(\overline B\).
Applying Lemma~1 to \(\resclass\) with its envelope \(\Cbar/2\) gives,
by \eqref{eq:complexity} and the sample-size assumption,
\[
 R_N(\resclass)+\tfrac{\Cbar}{2}\sqrt{2\ln(1/\delta)/N}
 \le \frac{\Cbar}{\sqrt N}\bigl(1+\tfrac12\sqrt{2\ln(1/\delta)}\bigr)
 \le \frac{\Xi(C,B,K,\delta)}{\sqrt N}\le\epsilon .
\]
Thus \(Z\ge\epsilon\) implies the event in \eqref{eq:lemma}, even at
equality. If \(\Cbar=0\), \eqref{eq:residual} forces every residual to vanish,
and \(\Prob\{Z\ge\epsilon\}=0\) since \(\epsilon>0\).
This proves \eqref{eq:theorem} with the printed threshold and confidence.

{\small\textbf{Sources.}
\href{https://moallemi.com/ciamac/papers/kernel-adp-2012.pdf\#page=29}{Bhat, Farias, Moallemi and Zheng, Appendix~B.1 and Appendix~C}.
\href{https://github.com/nikhilpb/kadp/blob/main/lean/Theorem3.lean}{Formal Theorem~3},
\href{https://github.com/nikhilpb/kadp/blob/main/lean/OneSidedConcentration.lean}{Lemma~1 and the zero-envelope example}, and
\href{https://github.com/nikhilpb/kadp/blob/main/lean/ResidualComplexity.lean}{residual complexity}.}
\end{document}
